%% kapa: exam=1 round=24 year=2013 period=2 form=A subject=appraisal_law subject_ko=부동산관계법규 seq=18 number=58
%% source: Q-Net 감정평가사 기출문제 (cst00309) artlSeq=5073802 file=extracted/1008/5073802/2013년 제24회 감정평가사 1차시험 기출문제/2교시(A형).hwp
%% conversion: kordoc 4.24.51 -> tools/gen1.py (auto); review=proofread-sonnet; flags=none
\begin{question}{58}
부동산 가격공시 및 감정평가에 관한 법령상 단독주택가격의 공시에 관한 설명으로 옳지 않은 것은?
\begin{choices}
  \item 국토교통부장관은 용도지역, 건물구조 등이 일반적으로 유사하다고 인정되는 일단의 단독주택 중에서 표준주택을 선정한다.
  \item 표준주택의 선정은 표준주택의 선정 및 관리지침에 따라야 한다.
  \item 국토교통부장관은 표준주택가격을 공시할 때에는 표준주택가격의 열람방법을 관보에 공고하여야 한다.
  \item 국토교통부장관이 따로 정하지 아니한 경우 표준주택가격의 공시기준일은 1월 1일로 한다.
  \item 국토교통부장관은 공시기준일 이후에 토지의 분할ㆍ합병이나 건물의 신축 등이 발생한 경우에는 대통령령이 정하는 날을 기준으로 하여 표준주택가격을 결정ㆍ공시하여야 한다.
\end{choices}
\end{question}
\answer{5}
